\documentclass{article}
\usepackage{amsmath}
\usepackage{amssymb}
\usepackage{booktabs}
\usepackage{geometry}
\usepackage{array}
\geometry{margin=1in}

\title{Pipeline Mathematics}
\author{}
\date{}

\begin{document}
\maketitle

% ─────────────────────────────────────────────────────────────────────────────
\section*{Step 1. Time Feature Augmentation}

$\mathbf{X} \in \mathbb{R}^{N \times 31}$: raw normalised input,
$N = 14{,}976$ province-week rows, 31 epidemiological, climate, and spatial features.
Missing values imputed with column means computed on training years only.

For each row $i$ with week index $w_i \in \{1,\ldots,52\}$:

\begin{align}
    t^{\sin}_{i} &= \sin\!\left(\tfrac{2\pi w_i}{52}\right), \qquad
    t^{\cos}_{i}  = \cos\!\left(\tfrac{2\pi w_i}{52}\right)
\end{align}

\begin{equation}
    \tilde{\mathbf{X}}
    = \bigl[\,\mathbf{X} \mid \mathbf{t}^{\sin} \mid \mathbf{t}^{\cos}\,\bigr]
    \in \mathbb{R}^{N \times 33}
\end{equation}

% ─────────────────────────────────────────────────────────────────────────────
\section*{Step 2. Autoencoder Encoding}

The encoder half of a trained autoencoder compresses each row:

\begin{equation}
    \tilde{\mathbf{X}} \in \mathbb{R}^{N \times 33}
    \;\xrightarrow{\text{encoder}}\;
    \mathbf{Z} \in \mathbb{R}^{N \times 6}
\end{equation}

$\mathbf{Z}$ contains the 6-dimensional latent representation of each province-week.

% ─────────────────────────────────────────────────────────────────────────────
\section*{Step 3. State Vector Construction}

Raw incidence $\mathbf{y} \in \mathbb{R}^{N}$ is clipped and log-transformed,
then appended as a 7th column:

\begin{equation}
    \tilde{y}_i = \log\!\bigl(1 + \max(0,\, y_i)\bigr)
\end{equation}

\begin{equation}
    \mathbf{S}
    = \bigl[\,\mathbf{Z} \mid \tilde{\mathbf{y}}\,\bigr]
    \in \mathbb{R}^{N \times 7},
    \qquad
    \mathbf{s}_i = \bigl[z_1,\, z_2,\, z_3,\, z_4,\, z_5,\, z_6,\, \tilde{y}\bigr]_i
\end{equation}

% ─────────────────────────────────────────────────────────────────────────────
\section*{Step 4. Lag Matrix Construction}

$p = \texttt{n\_lags}$ is a hyperparameter discovered by the island search (bounded in $[2,8]$, seeded at 5).

For every valid time point $t$ within a province, flatten the $p+1$ consecutive state rows into one wide feature vector:

\begin{equation}
    \boldsymbol{\theta}_t
    = \bigl[\,\mathbf{s}_{t-p} \mid \mathbf{s}_{t-p+1} \mid \cdots \mid \mathbf{s}_{t}\,\bigr]
    \in \mathbb{R}^{7(p+1)}
\end{equation}

The prediction target is the full state one week ahead:

\begin{equation}
    \mathbf{y}_t = \mathbf{s}_{t+1} \in \mathbb{R}^{7}
\end{equation}

Stacking all valid windows across all provinces:

\begin{equation}
    \boldsymbol{\Theta} \in \mathbb{R}^{T \times 7(p+1)},
    \qquad
    \mathbf{Y} \in \mathbb{R}^{T \times 7},
    \qquad T \approx 12{,}000
\end{equation}

For $p = 4$: $\boldsymbol{\Theta} \in \mathbb{R}^{T \times 35}$.
The first $p$ rows of each province chunk are dropped (insufficient history).

\medskip
\noindent\textbf{Lag interpretation.}
Each of the $p+1$ blocks in a row of $\boldsymbol{\Theta}$ is a different week,
so feature $z_k$ appears $p+1$ times as independent columns.
STLSQ fits a separate coefficient to each. After thresholding, only some survive.
If $z_k$ at lag $t-2$ survives but $z_k$ at $t$ and $t-1$ die,
the model has discovered that $z_k$ from two weeks ago drives the dynamics.
The lag structure is not imposed. It emerges from thresholding.

% ─────────────────────────────────────────────────────────────────────────────
\section*{Step 5. STLSQ}

\textbf{Objective.} Find a sparse $\boldsymbol{\Xi}$ such that:

\begin{equation}
    \mathbf{Y} \approx \boldsymbol{\Theta}\,\boldsymbol{\Xi},
    \qquad
    \boldsymbol{\Theta} \in \mathbb{R}^{T \times 35},\quad
    \boldsymbol{\Xi} \in \mathbb{R}^{35 \times 7},\quad
    \mathbf{Y} \in \mathbb{R}^{T \times 7}
\end{equation}

\textbf{OLS closed form} (one iteration):

\begin{equation}
    \hat{\boldsymbol{\Xi}}
    = \bigl(\boldsymbol{\Theta}^{\top}\boldsymbol{\Theta}\bigr)^{-1}
      \boldsymbol{\Theta}^{\top}\mathbf{Y}
\end{equation}

\begin{center}
\begin{tabular}{ll}
\toprule
Term & Shape \\
\midrule
$\boldsymbol{\Theta}^{\top}\boldsymbol{\Theta}$ & $35 \times 35$ \\
$\boldsymbol{\Theta}^{\top}\mathbf{Y}$          & $35 \times 7$  \\
$\hat{\boldsymbol{\Xi}}$                         & $35 \times 7$  \\
\bottomrule
\end{tabular}
\end{center}

\textbf{Iterative thresholding loop:}

\begin{align}
    &\text{1. Compute } \hat{\boldsymbol{\Xi}} \text{ via OLS.} \notag\\
    &\text{2. Zero out all entries: } \hat{\xi}_{ij} \leftarrow 0 \text{ if } |\hat{\xi}_{ij}| < \lambda. \notag\\
    &\text{3. Refit OLS on surviving columns only.} \notag\\
    &\text{4. Repeat until no further entries are eliminated.} \notag
\end{align}

The final sparse $\boldsymbol{\Xi}$ is the \textbf{discovered equation of motion}.
Its surviving entries are the coefficients of the governing equation:

\begin{equation}
    \mathbf{s}_{t+1} = \boldsymbol{\theta}_t\,\boldsymbol{\Xi}
\end{equation}

$\boldsymbol{\Xi}$ has one column per output state dimension
$[z_1, z_2, z_3, z_4, z_5, z_6, \tilde{y}]$ and one row per input feature
across all $p+1$ lag windows. Most entries are zero after thresholding.

% ─────────────────────────────────────────────────────────────────────────────
\section*{Step 6. Island Search over $(\lambda,\, p)$}

STLSQ requires two hyperparameters: threshold $\lambda$ and lag order $p$.
These are discovered by an island-model evolutionary search.

\subsection*{Island configuration}

\begin{center}
\begin{tabular}{ll}
\toprule
Parameter & Value \\
\midrule
Number of islands $K$       & 4 \\
Population per island       & 5 programs (best 5 retained) \\
Total iterations $M$        & 80 \\
Migration interval $G$      & every 10 iterations \\
Number of migration events  & $M / G = 8$ \\
\bottomrule
\end{tabular}
\end{center}

Each island is seeded with a different bias to explore different regions of the search space:

\begin{center}
\begin{tabular}{cll}
\toprule
Island & $p$ range & Bias \\
\midrule
0 & $[3,4]$   & conservative, short lags \\
1 & $[5,5]$   & default \\
2 & $[6,8]$   & deep-lag exploration \\
3 & $[4,5]$   & nonlinear, extra terms seeded \\
\bottomrule
\end{tabular}
\end{center}

\subsection*{One iteration within an island}

\begin{align}
    &\text{1. Sample parent } (\lambda, p) \text{ from island population via tournament selection.} \notag\\
    &\text{2. Mutate to produce child } (\lambda', p'). \notag\\
    &\text{3. Run STLSQ with } (\lambda', p') \text{ on training data.} \notag\\
    &\text{4. Score child on validation loss: } f = -\text{MSE}_{\text{val}}. \notag\\
    &\text{5. Add child to island; evict lowest-scoring program if population} > 5. \notag
\end{align}

\subsection*{Migration (every $G = 10$ iterations)}

Every 10 iterations, each island broadcasts its current best program to all other islands:

\begin{equation}
    \forall\, i \neq j : \quad
    \text{Island}_j.\text{add}\!\bigl(\text{Island}_i.\text{best}()\bigr)
\end{equation}

This is all-to-all migration. Migrated programs compete normally in the receiving island.
The full search timeline:

\begin{equation}
    \underbrace{[\,\text{10 iters} \to \text{migrate}\,]}_{\times\, 8}
    \quad\Rightarrow\quad
    \text{80 total STLSQ evaluations per island,}\quad
    8 \text{ migration events.}
\end{equation}

\subsection*{Final selection}

After all 80 iterations, each island reports its best $(\lambda^*, p^*)$ and corresponding $\boldsymbol{\Xi}^*$.
The winner is selected by validation score across all $K = 4$ islands:

\begin{equation}
    (\lambda^*, p^*) = \arg\max_{k \in \{1,\ldots,K\}} f_k^{\text{best}}
\end{equation}

The winning sparse $\boldsymbol{\Xi}^*$ and its governing equation are passed to EpiHSD for
evaluation on the held-out test set.

% ─────────────────────────────────────────────────────────────────────────────
\section*{Step 7. EpiHSD Input}

After the island search, the winning $(\lambda^*, p^*)$ and the corresponding sparse
$\boldsymbol{\Xi}^*$ define the governing equation.
EpiHSD then takes the same latent state vector $\mathbf{S}$ and runs a three-level
hierarchical sparse regression on it.

\subsection*{State vector entering EpiHSD}

The state vector is unchanged from Step 3:

\begin{equation}
    \mathbf{s}_t = [z_1, z_2, z_3, z_4, z_5, z_6, \tilde{y}]_t \in \mathbb{R}^7
\end{equation}

with $\tilde{y}_t = \log(1 + \max(0, y_t))$ and $L = 5$ lag steps.

\subsection*{Lag feature matrix for EpiHSD}

EpiHSD concatenates $L+1 = 6$ consecutive state blocks into one row
(same construction as Step 4, but with a fixed $p = L = 5$):

\begin{equation}
    \mathbf{x}_t
    = [\,\mathbf{s}_{t-L} \mid \cdots \mid \mathbf{s}_t\,]
    \in \mathbb{R}^{7 \times 6 = 42}
\end{equation}

A polynomial feature map of degree 1 (linear, with bias) is then applied:

\begin{equation}
    \boldsymbol{\theta}_t = \phi(\mathbf{x}_t) \in \mathbb{R}^{43}
    \qquad (\text{42 features} + \text{1 bias term})
\end{equation}

Stacking all valid windows across all provinces:

\begin{equation}
    \boldsymbol{\Theta} \in \mathbb{R}^{T \times 43},
    \qquad
    \mathbf{Y} \in \mathbb{R}^{T \times 7},
    \qquad T \approx 8{,}000 \text{ (train years)}
\end{equation}

\subsection*{Three-level hierarchical regression}

EpiHSD fits STLSQ in three nested residual levels.
Let $\mathbf{W} \in \mathbb{R}^{43 \times 7}$ denote a coefficient matrix at any level.

\medskip
\noindent\textbf{Level 1 — Global.}
Fit one equation on all provinces simultaneously:

\begin{equation}
    \mathbf{Y} \approx \boldsymbol{\Theta}\,\mathbf{W}_{\text{global}}
\end{equation}

\noindent\textbf{Level 2 — Cluster.}
For each cluster $c$, compute residuals and fit a sparse correction:

\begin{equation}
    \mathbf{R}^{(c)} = \mathbf{Y}^{(c)} - \boldsymbol{\Theta}^{(c)}\,\mathbf{W}_{\text{global}},
    \qquad
    \mathbf{R}^{(c)} \approx \boldsymbol{\Theta}^{(c)}\,\Delta\mathbf{W}^{(c)}
\end{equation}

\begin{equation}
    \mathbf{W}^{(c)} = \mathbf{W}_{\text{global}} + \Delta\mathbf{W}^{(c)}
\end{equation}

\noindent\textbf{Level 3 — Province.}
For each province $p$ in cluster $c$, fit a further sparse correction on top of the cluster equation:

\begin{equation}
    \mathbf{R}^{(p)} = \mathbf{Y}^{(p)} - \boldsymbol{\Theta}^{(p)}\,\mathbf{W}^{(c)},
    \qquad
    \mathbf{R}^{(p)} \approx \boldsymbol{\Theta}^{(p)}\,\Delta\mathbf{W}^{(p)}
\end{equation}

\begin{equation}
    \mathbf{W}^{(p)} = \mathbf{W}^{(c)} + \Delta\mathbf{W}^{(p)}
\end{equation}

\noindent The province-level equation $\mathbf{W}^{(p)}$ is the final governing equation for province $p$.

\subsection*{Forecasting}

Given the last $L+1$ weeks of observed state, one-step prediction for province $p$ is:

\begin{equation}
    \hat{\mathbf{s}}_{t+1} = \phi(\mathbf{x}_t)\,\mathbf{W}^{(p)} \in \mathbb{R}^7
\end{equation}

The 7th entry $\hat{s}_{t+1,7}$ is the predicted log-incidence.
Reporting converts back to raw scale: $\hat{y}_{t+1} = e^{\hat{s}_{t+1,7}} - 1$.

\subsection*{STLSQ thresholds at each level}

\begin{center}
\begin{tabular}{lcc}
\toprule
Level & Latent dims $z_1\ldots z_6$ & Incidence $\tilde{y}$ \\
\midrule
Global   & $\lambda = 0.05$ & $\lambda = 0.04$ \\
Cluster  & $\lambda = 0.05$ & $\lambda = 0.04$ \\
Province & $\lambda = 0.08$ & $\lambda = 100$ (disabled) \\
\bottomrule
\end{tabular}
\end{center}

The incidence threshold at province level is set to 100 (effectively infinity) so that
province-level corrections refine only the latent dynamics $z_1\ldots z_6$,
leaving incidence prediction entirely to the cluster equation.

% ─────────────────────────────────────────────────────────────────────────────
\section*{Step 8. Final Output: Discovered Governing Equations}

The full pipeline produces three nested coefficient matrices,
one per hierarchical level. Each is a sparse $\mathbf{W} \in \mathbb{R}^{43 \times 7}$
whose surviving entries are the discovered equation of motion.

\subsection*{How to read an equation}

Every surviving term has the form:

\begin{equation}
    \alpha \cdot z_{k,\, t - \ell}
    \quad \text{or} \quad
    \alpha \cdot \tilde{y}_{t-\ell}
\end{equation}

where $\alpha$ is the coefficient, $k \in \{1,\ldots,6\}$ is the latent dimension,
and $\ell \in \{0,1,\ldots,5\}$ is the lag. The output is one of the 7 state dimensions
at $t+1$.

\subsection*{Level 1: Global equation (example — incidence column)}

The globally discovered incidence equation is:

\begin{equation}
    \tilde{y}_{t+1}
    = 0.436\,\tilde{y}_{t}
    + 0.205\,\tilde{y}_{t-1}
    + 0.132\,\tilde{y}_{t-2}
    + 0.060\,\tilde{y}_{t-3}
    + 0.052\,\tilde{y}_{t-5}
    - 0.053\,z_{4,t}
    - 0.050\,z_{6,t}
\end{equation}

This says that across all 32 provinces, next-week log-incidence is driven primarily
by its own autoregressive history (lags 0-3 and 5), with small corrections from
latent dimensions $z_4$ and $z_6$ at the current week.

\subsection*{Level 2: Cluster delta equations}

Each cluster $c$ adds a sparse correction $\Delta\mathbf{W}^{(c)}$ on top of the global equation.
For example, Cluster 3 (11 provinces) has a minimal incidence correction:

\begin{equation}
    \Delta\tilde{y}^{(3)}_{t+1}
    = 0.055\,\tilde{y}_{t}
    + 0.060\,\tilde{y}_{t-1}
    - 0.116\,\tilde{y}_{t-2}
    + 0.050\,z_{4,t-3}
    + 0.045\,z_{1,t-1}
    + 0.062\,z_{4,t-1}
\end{equation}

Cluster 4 (the best-performing cluster, test $R^2 = 0.53$) has an even sparser correction,
relying almost entirely on the global autoregressive structure:

\begin{equation}
    \Delta\tilde{y}^{(4)}_{t+1}
    = 0.067\,\tilde{y}_{t}
    - 0.080\,\tilde{y}_{t-3}
    + 0.115\,\tilde{y}_{t-4}
    - 0.077\,\tilde{y}_{t-5}
\end{equation}

The sparsity of cluster corrections is by design. Clusters with dense $\Delta\mathbf{W}^{(c)}$
(e.g. Cluster 0 with 103 active delta terms) indicate that the global equation is a poor
fit for those provinces and more regional structure is needed.

\subsection*{Level 3: Province equation}

The final province-level equation is the sum of all three levels:

\begin{equation}
    \mathbf{W}^{(p)}
    = \mathbf{W}_{\text{global}}
    + \Delta\mathbf{W}^{(c(p))}
    + \Delta\mathbf{W}^{(p)}
\end{equation}

Province-level corrections refine only the latent dimensions $z_1\ldots z_6$.
Incidence prediction is frozen at the cluster level ($\lambda_{\text{inc}}^{(p)} = 100$).

\subsection*{Performance summary}

\begin{center}
\begin{tabular}{lcccc}
\toprule
Method & Val $R^2$ (inc) & Test $R^2$ (inc) & Interpretable & Hierarchical \\
\midrule
Linear Regression    & ---   & 0.611 & partial & no \\
Lasso                & ---   & 0.689 & partial & no \\
Random Forest        & ---   & 0.666 & no      & no \\
Gradient Boosting    & ---   & 0.665 & no      & no \\
GRU (neural)         & 0.160 & 0.290 & no      & no \\
EpiHSD baseline      & 0.230 & 0.397 & yes     & yes \\
EpiHSD + FunSearch   & 0.496 & 0.394 & yes     & yes \\
EpiHSD + FS cluster 4 & 0.691 & 0.531 & yes    & yes \\
\bottomrule
\end{tabular}
\end{center}

The simple ML baselines achieve higher raw incidence $R^2$ by fitting directly
to incidence history without any physical structure.
EpiHSD operates in a 6-dimensional latent space that encodes climate, spatial,
and epidemiological dynamics jointly. Its advantage is not peak $R^2$ on a single split
but the interpretability of the discovered equations, the hierarchical decomposition
of global versus regional dynamics, and the generalisable structure of the latent state.

\end{document}


